\section{Objective-cutoff propagation and its representation}
\label{sec:propagation}

An objective cutoff can certify a box before its relaxation is solved. Its
strength depends on the expression graph, and its cost depends on the
propagation schedule. We first separate these two dependencies. Throughout
this section, the graph contains the objective alone. Revising original
constraints jointly with the objective, imposing new auxiliary-variable
bounds, and branching on auxiliary variables define different schemes.
Natural interval extensions and validated exclusion regions are established
tools \citep{schichlMarkotNeumaier2014ExclusionRegions}. The results below
concern objective-cutoff propagation and the certificates it can produce
under a fixed graph.

Let $\mathcal D$ be a finite directed acyclic graph representing a continuous
function $f$ on a compact box $D$. Its nodes are variables, fixed constants,
and continuous operations on the compact forward domains. Each operation
has an elementary graph constraint $E$. Repeated references in such a
constraint refer to the same coordinate. For $C\subseteq D$, let $Z_0(C)$
be the lifted box obtained by forward interval evaluation, and let
$F_{\rm lo}(C)$ be its root lower endpoint. Forward evaluation is assumed
inclusion isotone and to use the exact image of each elementary operation
over its input intervals. This permits dependency loss between operations.

For a lifted box $Z$, the exact revise $\rho_E(Z)$ replaces the coordinates
in $E$ by the interval hull of their projections from $E\cap Z$; it returns
the empty box if this intersection is empty. Other coordinates remain
unchanged. The closed cutoff revise intersects the root interval with
$(-\infty,c]$. A box is \emph{hull-consistent} if it is fixed by every
elementary revise and by this cutoff revise. Hull consistency asks that
interval endpoints have support within each elementary constraint. Those
supports need not agree between constraints.
An exact schedule is fair when every elementary revise and the cutoff
revise occur infinitely often. For partial revises, fairness requires both
the forward and backward revise of every operation, and the cutoff revise,
infinitely often.

\begin{theorem}[The objective-DAG bound]
\label{prop:fixed-point}
There is a greatest hull-consistent box $Z^*(C,c)\subseteq Z_0(C)$.
Every sequence of steps satisfying
$\rho_E(Z)\subseteq\rho(Z)\subseteq Z$ preserves this box. A fair sequence
of exact revises converges to $Z^*(C,c)$, and reaches an empty iterate if
this limit is empty. The same conclusions hold for fair exact forward and
backward partial revises.

For $y\in C$, the minima
\begin{align}
 \pi_{\mathcal D}(C)&=\min\{c:Z^*(C,c)\ne\varnothing\},\label{prop:pi-box}\\
 \pi_{\mathcal D}(C,y)&=\min\{c:y\in\Pi_xZ^*(C,c)\}\label{prop:pi-point}
\end{align}
exist and satisfy
\begin{equation}
 F_{\rm lo}(C)\le\pi_{\mathcal D}(C)
 \le\pi_{\mathcal D}(C,y)\le f(y).
 \label{prop:bound-chain}
\end{equation}
Both bounds increase when $C$ shrinks. The limiting propagation empties $C$
exactly when $c<\pi_{\mathcal D}(C)$ and removes $y$ exactly when
$c<\pi_{\mathcal D}(C,y)$. Any allowed finite schedule removes only points
satisfying this strict inequality.
\end{theorem}

The fixed-point viewpoint is standard in FBBT
\citep{belottiCafieriLeeLiberti2012FBBT}. The proof for the present schedule
and cutoff model is in \cref{app:propagation-fixed-point}. The strict comparisons
matter: a closed cutoff remains consistent at its threshold. Relaxation
pruning, in contrast, can use $L(C)\ge f_* -\epsilon$. Computing the limiting
box as a node oracle does not provide a finite stopping test for every
nonempty limit.

\begin{definition}[Flat objective sum]
\label{prop:flat-definition}
The root equation is
$w=b+\sum_{j=1}^J a_jp_j$, with nonzero coefficients and \emph{distinct}
root-child node coordinates $p_j$. Each $p_j$ is a base node $z_{k(j)}$ or
a unary operation $p_j=\varphi_j(z_{k(j)})$ used only by the root. Write
$t_j=a_j\varphi_j$, with $\varphi_j(z)=z$ for direct base terms. A base can
feed several distinct unary nodes. Repeated references to the same root
coordinate must first be collected into one coefficient. For base intervals
$U_k=[s_k,t_k]$, set
\begin{align}
 \mu_j(U)&=\min_{v\in U_{k(j)}}t_j(v),&
 h_j(U)&=\max\{t_j(s_{k(j)}),t_j(t_{k(j)})\},\\
 \Phi(U)&=b+\sum_j\mu_j(U)+\max_j\{h_j(U)-\mu_j(U)\}.
 \label{prop:flat-functional}
\end{align}
\end{definition}

\begin{theorem}[Exact flat-sum formula]
\label{prop:flat-formula}
If $Z^*(C,c)$ is nonempty, its base intervals $U$ satisfy $\Phi(U)\le c$.
If the bases are the original variables, then
\begin{equation}
 \pi_{\mathcal D}(C)=\min_{\varnothing\ne U\subseteq C}\Phi(U),\qquad
 \pi_{\mathcal D}(C,y)=\min_{y\in U\subseteq C}\Phi(U),
 \label{prop:flat-minima}
\end{equation}
where $U$ ranges over closed subboxes, including degenerate ones.
\end{theorem}

\begin{proof}
In a consistent box, the signed term interval $T_j=a_jZ_{p_j}$ has lower
endpoint at least $\mu_j$, and upper endpoint at least $h_j$: each base
endpoint must have support in its unary constraint. Supporting the upper
endpoint of $T_j$ in the root constraint gives
$c\ge b+h_j+\sum_{i\ne j}\mu_i$. This proves the necessary inequality.

Conversely give the variables the intervals $U$, give each signed term
the interval $[\mu_j,h_j]$, and give the root
$[b+\sum_j\mu_j,\min\{c,b+\sum_jh_j\}]$. Unary constraints support their
base endpoints and their term endpoints. Direct variable terms have
their exact signed intervals. For the root, any term endpoint can be
supported by putting all other terms at their minima; the resulting sum
is at most $\Phi(U)\le c$. The root endpoints are attainable in a sum
of intervals. Distinct root coordinates make these root supports legitimate.
This constructs a consistent box containing $U$. The necessary inequality
applied to $\Pi_x Z^*$ proves both formulas. Continuity of term minima and
endpoint values, on the compact set of subbox endpoints, proves attainment.
\end{proof}

The endpoint maximum in \eqref{prop:flat-functional} cannot be replaced by
the full range maximum. For example, representing $x^2$ on $[-1,1]$ by
three distinct power nodes with coefficients $-3,2,2$ gives
$\Phi([-a,a])=-a^2$, and hence $\pi_{\mathcal D}([-1,1])=-1$.
Indeed a same-side subinterval gives $\Phi=\min_Ux^2\ge0$; a subinterval
crossing zero gives $\Phi\ge-\max_Ux^2\ge-1$.
The single-node square representation has bound zero. Moreover the formula
does not apply to the root equation $w=x+x-x$: its repeated references
give the exact equation $w=x$, not three independent root coordinates.

\begin{corollary}[Sufficient exactness in exposed coordinates]
\label{prop:one-sided}
For one base coordinate, call an interval $U$ one-sided if there are a term
$j_0$ and an endpoint $e$ such that
$t_{j_0}(e)=h_{j_0}(U)$ and $t_j(e)=\mu_j(U)$ for every $j\ne j_0$.
If every subinterval of its forward base range is one-sided, then
\[
 \pi_{\mathcal D}(C)\ge\min_{z\in Z_0(C)_z}
 \left(b+\sum_jt_j(z)\right).
\]
When the base is the original variable, this is equality with $\min_Cf$.

More generally write $\widehat f(z)=b+\sum_jt_j(z_{k(j)})$ and let
$B_z=\prod_kZ_0(C)_{z_k}$. For each coordinate group define
\[
 e_k(U_k)=\min_{U_k}\sum_{j:k(j)=k}t_j-
                   \sum_{j:k(j)=k}\mu_j\ge0.
\]
If on every subbox of $B_z$ each coordinate group is one-sided and at
most one $e_k$ is positive, then
$\pi_{\mathcal D}(C)\ge\min_{B_z}\widehat f$.
This bound is exact if $z(C)=B_z$.
\end{corollary}

\begin{proof}
In the first case $\Phi(U)\ge\widehat f(e)$.
For the second, put
$G_k=\max_{j:k(j)=k}(h_j-\mu_j)-e_k$. One-sidedness gives $G_k\ge0$, and
direct substitution gives
\[
 \Phi(U)=\min_U\widehat f+
 \max_k\left(G_k-\sum_{\ell\ne k}e_\ell\right).
\]
Choose the sole coordinate with positive excess, or any coordinate if
all excesses vanish. The last maximum is nonnegative. Apply the necessary
part of \cref{prop:flat-formula} and \eqref{prop:bound-chain}.
\end{proof}

Thus exposing a coordinate such as $z=x_1-x_2$ can change certification.
For example $h(z)=(z-1)^2((z+1)^2+1/2)$, written as
$z^4-1.5z^2-z+1.5$, is one-sided on every positive interval: only $z^4$
increases. Adding a second exposed coordinate with one nonnegative square
term preserves the sufficient criterion. These are sufficient conditions;
they do not classify all exact graphs.

\begin{theorem}[Local exactness and ideal node counts]
\label{prop:local-exactness}
Assume $F=D$, $f$ is $L$-Lipschitz in the sup norm, and
$L(C)\ge\min_Cf-\tau w(C)^2$, where $w(C)$ is the largest side length.
Suppose a relatively open set $R$ contains the minimizing set and, for
some $w_0>0$,
$\pi_{\mathcal D}(C)\ge f_*$ for all $C\subseteq R$ with $w(C)\le w_0$.
An ideal algorithm with incumbent $f_*$ first decides whether the objective
fixed point at cutoff $f_* -\epsilon$ is empty, then optionally contracts
to a box containing its variable projection, then applies the relaxation
test and widest-side bisection. Its node count is bounded independently of
$\epsilon>0$.
\end{theorem}

\begin{proof}
If $R=D$, boxes of width at most $w_0$ are pruned by propagation.
Otherwise $\delta=\min_{D\setminus R}(f-f_*)>0$. Choose $0<h\le w_0$ with
$Lh\le\delta/2$ and $\tau h^2\le\delta/2$. A box of width at most $h$
is either contained in $R$, or contains a point outside $R$. In the latter
case $f-f_*\ge\delta/2$ throughout it, so any contracted box passes the
relaxation test. Thus no such box is split. The integer potential
$\sum_i\max\{0,\lceil\log_2(w_i(C)/h)\rceil\}$ never increases under
contraction and falls by at least one at each split. If its root value is
$d$, the tree has at most $2^{d+1}-1$ nodes.
\end{proof}

This theorem concerns the fixed-point oracle. Its decision may require
many rounds. It is the exact-bound mechanism behind absence of clustering
\citep{duKearfott1994Cluster,wechsungSchaberBarton2014Cluster}, applied to
a graph-dependent bound. No representation-free node lower bound can hold
in this unconstrained model: the representation
$f_*+|f-f_*|$ equals $f$ on $D$ and has forward lower bound $f_*$.
It encodes the answer in its constant $f_*$. For a constrained optimum this
identity need not hold at infeasible points.

\begin{lemma}[A surviving-point witness]
\label{prop:witness}
Suppose the root is a sum of distinct child nodes, over any continuous
DAG below them. Regard their signed expressions as $t_j(x)$. For every
$y\in U\subseteq C$,
\begin{equation}
 \pi_{\mathcal D}(C,y)\le
 b+\sum_j\min_Ut_j+\max_j\left(\max_Ut_j-\min_Ut_j\right).
 \label{prop:full-witness}
\end{equation}
If each term is an expression tree with each variable used at most once,
then its forward range is exact and
\[
 F_{\rm lo}(C)\le\pi_{\mathcal D}(C)
 \le F_{\rm lo}(C)+\max_j\bigl(\max_Ct_j-\min_Ct_j\bigr).
\]
\end{lemma}

\begin{proof}
Give every nonroot node its exact range over $U$. Every value in each
interval is attained by a point of $U$, and therefore has support within
its elementary graph constraint, even if different values use different
points. Give the root the interval from the proof of
\cref{prop:flat-formula}, now using full term maxima. This gives a
consistent witness with variable box $U$. For a single-use term, the
children of each binary operation use disjoint variables, so their joint
range is the product of their ranges. Induction proves exact forward
evaluation, giving the second assertion.
\end{proof}

\begin{theorem}[Coordinatewise loss preserves geometric lower bounds]
\label{prop:cnd-transfer}
Assume $F=D$ and the root is a sum of distinct term nodes with $C^2$
signed expressions $t_j$. Let $s_0>0$ be the largest root side. On a set
$N\subseteq D$, assume, for every coordinate,
\[
 d_i(y):=\sum_j|\partial_it_j(y)|-2\max_j|\partial_it_j(y)|\ge d_0>0.
\]
On coordinate segments of radius $\rho>0$ about these points, suppose
$|\partial_{ii}t_j|\le H_j$, and put
$K_0=\sum_jH_j/2+\max_jH_j$,
$\rho_* =\min\{\rho,d_0/(2K_0)\}$,
$\vartheta=d_0\rho_*/2$, interpreting $d_0/(2K_0)=+\infty$ if $K_0=0$.
If propagation removes $y\in N\cap C$ at cutoff $f_* -\epsilon$, and
$m(y)+\epsilon<\vartheta$, then
\begin{equation}
 \min\{y_i-l_i,u_i-y_i\}<\frac{2(m(y)+\epsilon)}{d_0},\qquad
 m(y)+\epsilon\ge\frac{d_0}{2ns_0}q_C(y).
 \label{prop:cnd-validity}
\end{equation}

Consider finite runs using coordinate splits, objective-only propagation
phases with inherited lifted bounds, and a pointwise relaxation test that
certifies $m(y)+\epsilon\ge\alpha q_C(y)$ on each discarded owned region.
Let $K$ count terminal boxes plus at most $2n$ removed-piece certificates
per phase. Set $\alpha_{\rm eff}=\min\{\alpha,d_0/(2ns_0)\}$. Then
\begin{align}
 K&\ge 2^{-n}\mathcal N_\infty\left(E(\eta)\cap N,
                     \sqrt{(\epsilon+\eta)/\alpha_{\rm eff}}\right),
      &&\epsilon+\eta<\vartheta,\label{prop:cnd-cover}\\
 K&\ge\left(\frac{\alpha_{\rm eff}n}{\pi^2}\right)^{n/2}
       \int_{N\cap\{m+\epsilon<\vartheta\}}
             (m(y)+\epsilon)^{-n/2}\,dy.
 \label{prop:cnd-integral}
\end{align}
Here $E(\eta)=\{m\le\eta\}$ and $\mathcal N_\infty(S,r)$ counts
sup-norm balls of radius $r$ covering $S$.
\end{theorem}

The full event and inheritance proof is in
\cref{app:propagation-transfer}. In particular, shared closed boundaries
stay with the retained box; the discarded owner, rather than every point
of its closed test box, satisfies the strict propagation comparison.
A $p$-dimensional smooth piece of minimizers in $N$ therefore forces
$K=\Omega(\epsilon^{-p/2})$. An interior minimizer with a neighborhood
contained in $N$ and
$m(y)\le\Gamma\|y-x_*\|^2$ locally forces
$K=\Omega(\log(1/\epsilon))$. These become node lower bounds when the
number of phases per node is uniformly bounded. Arbitrarily many phases
cannot be charged to a bounded number of nodes.

Coordinatewise no-dominance is stronger than the aggregate condition
\[
 \ell(y)=\sum_j\|\nabla t_j(y)\|_1
                  -2\max_j\|\nabla t_j(y)\|_1>0.
\]
Aggregate loss confines removed points to thin face layers, rather than
to vertex neighborhoods. Under uniform aggregate loss, an interior
nondegenerate minimum in dimension $n\ge2$ still forces a logarithmic
event count; a minimizing arc whose tangent has every coordinate bounded
away from zero still forces $\Omega(\epsilon^{-1/2})$ events. Precise
hypotheses and complete proofs are in \cref{app:propagation-face-loss}.
These weaker hypotheses do not establish the general
$\epsilon^{-p/2}$ rate in dimension $p\ge3$.

Finally, the node reduction in \cref{prop:local-exactness} can move work
into propagation rounds. For the graph $s^2-2as+a^2$ on a fixed interval
around $a>0$, a forward-cutoff-backward schedule needs
$\Omega(a\epsilon^{-1/2})$ rounds before emptying, although its limiting
bound is exact. For $u=t^2$, $v=u^2$, $f=u-2v$, a schedule with initial
$0<U_0<1/2$ on $u$ empties in $O(\log\log(1/\epsilon))$ rounds.
\Cref{app:propagation-rounds} proves both statements for the specified
schedules. Neither is a bound for all graphs or all schedules. General
criteria for bounded-round propagation remain open.
